% Supplementary Material for "GPU-Scale Experiments on Random 3-SAT"
% Author: Nadim Faris Nadim Zaro

\documentclass[11pt]{article}
\usepackage[utf8]{inputenc}
\usepackage[margin=1in]{geometry}
\usepackage{amsmath,amsthm,amssymb}
\usepackage{graphicx}
\usepackage{hyperref}
\usepackage{booktabs}
\usepackage{listings}
\usepackage{xcolor}
\usepackage{multirow}

\title{\textbf{Supplementary Material:} \\ GPU-Scale Experiments on Random 3-SAT}
\author{
  Nadim Faris Nadim Zaro \\
  \textit{Independent Researcher} \\
  \texttt{zaronadim@gmail.com}
}
\date{September 3, 2026}

% Define custom commands
\newcommand{\SAT}{\textsf{SAT}}
\newcommand{\NP}{\textsf{NP}}
\newcommand{\PTIME}{\textsf{P}}

\begin{document}

\maketitle

\tableofcontents

\section{Extended Methodology}

\subsection{Detailed GPU Implementation}

\textbf{Hardware Specifications:}
\begin{itemize}
\item GPU: NVIDIA GeForce RTX 5070 Laptop GPU
\item Compute Capability: 12.0
\item Multiprocessors: 36
\item CUDA Cores: 4,608 (estimated)
\item Memory: 8.5 GB GDDR7
\item Memory Bandwidth: ~448 GB/s
\item Peak FP32 Performance: ~15 TFLOPS
\end{itemize}

\textbf{Software Stack:}
\begin{itemize}
\item CUDA: 12.x
\item CuPy: 13.x
\item Python: 3.13
\item NumPy: 1.26
\item Operating System: Windows 11
\end{itemize}

\subsection{SAT Enumerator Optimization Details}

\textbf{Memory Management:}
\begin{lstlisting}[language=Python]
# Chunk size selection based on GPU memory
if n_vars <= 26:
    chunk_size = 2**26  # 67M assignments
elif n_vars <= 28:
    chunk_size = 2**25  # 33.5M assignments
else:
    chunk_size = 2**24  # 16.8M assignments
\end{lstlisting}

The reduced chunk size for $n \geq 29$, together with increased memory
pressure, is the most plausible cause of the throughput drop observed at
$n=30$ in the main paper (Experiment 1).

\textbf{Parallelization Strategy:}
\begin{itemize}
\item Each GPU thread evaluates one clause across all assignments
\item Clause results combined via parallel AND reduction
\item Memory coalescing for optimal bandwidth
\item Shared memory for clause data
\end{itemize}

\textbf{Throughput Optimization:}
\begin{enumerate}
\item Bit-manipulation for assignment generation
\item Vectorized clause evaluation
\item Early termination on first solution found (this makes per-instance
      ``throughput'' depend on where the first solution happens to lie)
\item GPU memory pooling to avoid allocation overhead
\end{enumerate}

\subsection{Sampled Solution-Graph Statistics}

For each instance, up to 400--500 assignments are drawn uniformly at random;
those satisfying the formula are retained. Distances are computed on GPU:
\begin{lstlisting}[language=Python]
solutions_gpu = cp.array(solutions)  # Shape: (N, n)
diff = solutions_gpu[:, cp.newaxis, :] -
       solutions_gpu[cp.newaxis, :, :]
distances = cp.sum(cp.abs(diff), axis=2)
\end{lstlisting}

A graph is built at the \emph{median} pairwise distance threshold:
\[
L = D - A \quad \text{where} \quad A[i,j] = \mathbb{1}[d(i,j) \leq t],
\]
and we report
\begin{align}
\beta_0 &= \#\{\lambda(L) : \lambda \approx 0\} \quad \text{(connected components)} \\
\beta_1 &= |E| - |V| + \beta_0 \quad \text{(graph cycle rank)}.
\end{align}

Interpretation notes: this is the cycle rank of one graph at one
data-dependent threshold (no filtration is computed, so it is not persistent
homology). Because the threshold is the median distance, about half of all
vertex pairs are edges by construction, so $\beta_1$ grows roughly
quadratically with the number of retained solutions --- which is governed by
clause density. The large cross-density differences in $\beta_1$ reported
below are therefore primarily a solution-density and sample-size effect.

\subsection{Instance Feature Extraction}

128-dimensional feature vector $\phi(I)$ for SAT instance $I$:

\textbf{Basic Features (16 dimensions):}
\begin{itemize}
\item $n$ (variables)
\item $m$ (clauses)
\item $m/n$ (clause-to-variable ratio)
\item Clause length statistics (mean, std, min, max)
\item Variable degree statistics
\end{itemize}

\textbf{Graph Features (32 dimensions):}
\begin{itemize}
\item Variable graph degree centrality
\item Clause graph clustering coefficient
\item Community structure metrics
\item Graph diameter and radius
\end{itemize}

\textbf{Structural Hashes (80 dimensions):}
\begin{itemize}
\item Hash of clause patterns
\item Symmetry group size
\item Backbone variables (forced assignments)
\item Resolution complexity estimates
\end{itemize}

The pairwise quantity reported in the main paper is
$R(P_1,P_2) = \|\phi(P_1) - \phi(P_2)\|_2$: a dissimilarity between instance
encodings, used as a descriptive statistic of the instance distribution.

\section{Complete Experimental Results}

\subsection{Full SAT Enumeration Table}

\begin{table}[h]
\centering
\caption{Complete Extreme-Scale SAT Results}
\begin{tabular}{ccccccc}
\toprule
$n$ & Assignments & Time (s) & Throughput (M/s) & SAT? & Solutions & Density \\
\midrule
26 & 67,108,864 & 2.05 & 32.7 & No & 0 & 4.2n \\
27 & 134,217,728 & 1.06 & 126.2 & Yes & 1 & 3.0n \\
28 & 268,435,456 & 1.07 & 250.5 & Yes & 8 & 2.0n \\
29 & 536,870,912 & 0.55 & 978.8 & Yes & 2 & 4.267n \\
30 & 1,073,741,824 & 7.61 & 141.1 & Yes & 1 & 4.2n \\
\midrule
\textbf{Total} & \textbf{2,080,374,784} & \textbf{12.34} & - & - & - & - \\
\bottomrule
\end{tabular}
\end{table}

Satisfiable instances terminate at the first solution found, so the time and
throughput columns mix hardware behavior with instance luck; the $n=29$ run
terminated after scanning roughly half its space, while the $n=30$ run
scanned most of it.

\textbf{Additional Statistics:}
\begin{itemize}
\item Average throughput: 168.6 M/s
\item Peak throughput: 978.8 M/s (n=29)
\item Satisfiability rate: 80\% (4/5)
\item GPU utilization: 87--94\%
\item Memory usage: 6.2--7.8 GB
\end{itemize}

\subsection{Solution-Graph Statistics --- Extended Results}

\begin{table}[h]
\centering
\caption{Cycle rank $\beta_1$ by clause density (3000 instances). Labels in
parentheses are the category tags used in the code and result files; all five
generators draw uniformly random 3-clauses and differ only in density.}
\begin{tabular}{lrrrrr}
\toprule
Density (code label) & $n$ & Mean $\beta_1$ & Std $\beta_1$ & Max $\beta_1$ & Instances \\
\midrule
2.0n (``algebraic'') & 600 & 126.72 & 411.24 & 5,512 & 600 \\
3.0n (``hierarchical'') & 600 & 8.92 & 40.38 & 443 & 600 \\
4.267n (``phase\_transition'') & 600 & 0.39 & 2.74 & 44 & 600 \\
5.0n (``planted'') & 600 & 0.35 & 4.19 & 82 & 600 \\
4.2n (``random'') & 600 & 1.26 & 7.49 & 83 & 600 \\
\bottomrule
\end{tabular}
\end{table}

The monotone trend --- $\beta_1$ collapses as density rises --- is what the
sampling procedure predicts: at density 2.0 many uniformly random assignments
satisfy the formula (large sampled graphs, quadratically many edges at the
median threshold), while at density $\geq 4.2$ almost none do (graphs with
$<3$ vertices are assigned $\beta_1 = 0$).

\textbf{Statistical Tests:}
\begin{itemize}
\item Kruskal-Wallis H-test: $H = 1847.2$, $p < 10^{-400}$
\item Post-hoc pairwise comparisons (Dunn's test):
  \begin{itemize}
  \item 2.0n vs. 4.2n: $p < 10^{-300}$
  \item 2.0n vs. 3.0n: $p < 10^{-200}$
  \item All other pairs: $p < 0.001$
  \end{itemize}
\end{itemize}

These tests confirm that clause density strongly affects the sampled
statistic.

\subsection{Feature-Distance Matrix --- Full Statistics}

\begin{table}[h]
\centering
\caption{Feature-space distance distribution}
\begin{tabular}{lr}
\toprule
Statistic & Value \\
\midrule
Instance pairs analyzed & 441,000 \\
Matrix shape & $441 \times 1000$ \\
Computation time & 1.17 seconds \\
Throughput & 377,218 pairs/sec \\
\midrule
Minimum & 4.0 \\
1st Quartile (Q1) & 18.2 \\
Median (Q2) & 25.0 \\
Mean & 31.0 \\
3rd Quartile (Q3) & 42.8 \\
Maximum & 110.5 \\
\midrule
Interquartile Range (IQR) & 24.6 \\
Standard Deviation & 18.7 \\
Coefficient of Variation & 0.60 \\
\bottomrule
\end{tabular}
\end{table}

\subsection{Neural Network --- Training Details}

\textbf{Task:} binary classification of an instance-size category: label 0 =
generated instances with $\leq 15$ variables (one dataset slice), label 1 =
instances of arbitrary size (another slice). All instances in both classes
are random 3-\SAT{}.

\textbf{Architecture:}
\begin{verbatim}
Layer 1: Dense(128 -> 512) + ReLU + Dropout(0.3)
Layer 2: Dense(512 -> 256) + ReLU + Dropout(0.3)
Layer 3: Dense(256 -> 128) + ReLU + Dropout(0.2)
Layer 4: Dense(128 -> 2) + Softmax
\end{verbatim}

\textbf{Training Hyperparameters:}
\begin{itemize}
\item Optimizer: Adam($\beta_1=0.9, \beta_2=0.999, \epsilon=10^{-8}$)
\item Learning rate: 0.001 (fixed)
\item Batch size: 512
\item Epochs: 150
\item Loss: Categorical cross-entropy
\item Weight initialization: He normal
\end{itemize}

\textbf{Training Progression:}
\begin{table}[h]
\centering
\caption{Neural Network Training Curve}
\begin{tabular}{ccc}
\toprule
Epoch & Training Loss & Best Loss \\
\midrule
0 & 1.4143 & 1.4143 \\
20 & 0.5440 & 0.5440 \\
40 & 0.5482 & 0.5414 \\
60 & 0.5389 & 0.5320 \\
80 & 0.5388 & 0.5320 \\
100 & 0.5361 & 0.5262 \\
120 & 0.5259 & 0.5259 \\
140 & 0.5267 & 0.5220 \\
\bottomrule
\end{tabular}
\end{table}

\textbf{Confusion Matrix (1000 test instances):}
\begin{table}[h]
\centering
\begin{tabular}{cc|cc}
& & \multicolumn{2}{c}{Predicted} \\
& & small & mixed \\
\hline
\multirow{2}{*}{Actual} & small & 356 & 144 \\
& mixed & 78 & 422 \\
\end{tabular}
\end{table}

\textbf{Performance Metrics:}
\begin{itemize}
\item Accuracy: 77.8\%
\item Precision: 0.79
\item Recall: 0.76
\item F1-Score: 0.77
\item ROC-AUC: 0.84
\end{itemize}

Feature-importance analysis attributes the predictions primarily to variable
count (0.83) and clause density (0.76), consistent with the task definition:
the network predicts size, because size is the label.

\section{Reproducibility}

\subsection{Code and Data Availability}

All code and data are publicly available:

\begin{itemize}
\item \textbf{GitHub Repository:} \url{https://github.com/big-brain-zaru/PvsNP}
\item \textbf{Zenodo Archive (DOI):} \url{https://doi.org/10.5281/zenodo.18451653}
\end{itemize}

\subsection{System Requirements}
\begin{itemize}
\item NVIDIA GPU with compute capability $\geq 7.0$
\item 8+ GB GPU memory
\item CUDA 11.0+
\item Python 3.9+
\item CuPy, NumPy, SciPy
\end{itemize}

\subsection{Running the Analysis}
\begin{lstlisting}[language=bash]
# Install dependencies
pip install cupy-cuda12x numpy scipy

# Clone repository
git clone https://github.com/big-brain-zaru/PvsNP
cd PvsNP

# Main analysis: generation, enumeration, solution-graph
# statistics, feature distances, classifier
python gpu_pnp_breakthrough.py

# Extreme-scale enumeration benchmarks (n = 26..30)
python gpu_formalization.py

# Aggregation of the JSON outputs
python ultimate_proof.py
\end{lstlisting}

Notes on script output: \texttt{gpu\_formalization.py} additionally prints an
``oracle relativization'' section whose values are produced by a hardcoded
simulation stub (fixed multipliers on a baseline constant), and
\texttt{ultimate\_proof.py} prints heuristic per-method verdict labels and an
aggregate score. Neither of these outputs is a measurement, and neither is
reported in the paper; the scripts are kept as-is so that the published JSON
result files remain exactly reproducible.

\textbf{Expected runtime:}
\begin{itemize}
\item Main analysis: 20--25 minutes
\item Extreme-scale benchmarks: 10--15 minutes
\item Aggregation: 2--3 minutes
\item Total: ~40 minutes on RTX 5070
\end{itemize}

\section{Conclusion}

This supplement documents the implementation and complete measurements behind
the main paper: the GPU enumeration pipeline and its memory management, the
sampled solution-graph statistics and their density dependence, the feature
extraction and distance statistics, and the classifier training details. The
measurements stand as benchmarking and instance-structure data for random
3-\SAT{}; their scope is discussed in the main paper.

\end{document}
